\documentclass[11pt,a4paper]{report}

\usepackage[utf8]{inputenc}
\usepackage[T1]{fontenc}
\usepackage[brazil]{babel}
\usepackage{lmodern}
\usepackage{amsmath,amssymb,amsthm}
\usepackage{booktabs}
\usepackage{longtable}
\usepackage{array}
\usepackage{xcolor}
\usepackage{geometry}
\usepackage{tikz}
\usepackage{pgfplots}
\usepackage{fancyhdr}
\usepackage{enumitem}
\usepackage{float}
\usepackage[hidelinks]{hyperref}

\pgfplotsset{compat=1.17}
\usetikzlibrary{arrows.meta,positioning,shapes.geometric,fit,backgrounds,calc}

\geometry{left=3cm,right=2.7cm,top=2.8cm,bottom=2.8cm}

\definecolor{rule}{RGB}{190,190,185}
\definecolor{callout}{RGB}{242,244,248}
\definecolor{accent}{RGB}{20,70,140}
\definecolor{warm}{RGB}{176,72,32}
\definecolor{soft}{RGB}{150,150,150}
\definecolor{good}{RGB}{30,120,70}

\newcommand{\cpn}{\textsc{CPNCheck}}
\newcommand{\pyra}{\textsc{pyRuleAnalyzer}}
\newcommand{\cpntools}{CPN~Tools}
\newcommand{\code}[1]{\texttt{\small #1}}

\newenvironment{nota}
  {\par\vspace{6pt}\noindent\begin{tikzpicture}\node[fill=callout,rounded corners=2pt,
   inner sep=9pt,text width=\dimexpr\linewidth-20pt\relax,align=justify]\bgroup\small}
  {\egroup;\end{tikzpicture}\par\vspace{6pt}}

\theoremstyle{definition}
\newtheorem{definicao}{Definição}[chapter]
\newtheorem{teorema}{Teorema}[chapter]
\newtheorem{exemplo}{Exemplo}[chapter]

\pagestyle{fancy}
\fancyhf{}
\fancyhead[L]{\small\nouppercase{\leftmark}}
\fancyhead[R]{\small\thepage}
\renewcommand{\headrulewidth}{0.4pt}

\title{\vspace{-1.5cm}\Huge\normalfont\scshape CPNCheck\upshape\\[6pt]
\Large Tratamento da explosão de estados\\[4pt]
\large redução por ordem parcial com conjuntos teimosos (\emph{stubborn sets})}
\author{Documentação técnica}
\date{\today}

\begin{document}

\maketitle
\thispagestyle{empty}

\vspace{1cm}
\begin{center}
\begin{minipage}{0.84\textwidth}
\small
\noindent\textbf{Sobre este documento.} Ele descreve a etapa do \cpn{} que
combate a explosão de estados: de onde a explosão vem nas redes geradas pelo
\pyra{}, quais técnicas da literatura foram avaliadas (inclusive as três do
capítulo~8 de Jensen \& Kristensen), por que a escolhida foi a redução por
ordem parcial com \emph{conjuntos teimosos}, como ela funciona passo a passo,
o que exatamente ela garante, onde ela para de funcionar, e como foi validada.

\vspace{6pt}
\noindent Não há código aqui. O texto explica o \emph{funcionamento} e o
\emph{porquê}; a implementação está em \code{cpncheck/reduction.py} e os
testes em \code{tests/test\_reduction.py}.
\end{minipage}
\end{center}

\vfill
\begin{center}
\small
\url{https://github.com/GrupoCybersegurancaVirtus/cpncheck}
\end{center}

\clearpage
\tableofcontents
\clearpage

% ============================================================================
\chapter{O problema: de onde vem a explosão}
\label{cap:problema}
% ============================================================================

\section{O que o model checker enumera}

O \cpn{} decide propriedades temporais (``a rede sempre termina'', ``a
predição nunca recebe dois tokens'', ``toda execução chega a uma decisão'')
construindo o \emph{grafo de ocorrência}: um nó para cada marcação alcançável a
partir da marcação inicial, um arco para cada ocorrência de um elemento de
ligação (uma transição com valores para suas variáveis). As fórmulas CTL são
então avaliadas sobre esse grafo por pontos fixos.

O custo de tudo é o tamanho desse grafo. Se ele tem milhões de nós, cada
propriedade custa milhões de passos; se tem $10^{30}$, nenhuma é decidida.
Isso é a \emph{explosão de estados}, a principal limitação dos métodos de
espaço de estados (Jensen \& Kristensen, 2009, cap.~8).

\section{O que foi medido nas redes do \pyra{}}

Antes de escolher uma técnica, foi medido \emph{o quê} explode. As redes
geradas a partir de modelos do scikit-learn têm uma estrutura fixa
(Figura~\ref{fig:rede-rf}): uma transição \code{Distribute} copia a amostra
para a entrada de cada árvore; cada árvore é uma página cujas transições são as
folhas, uma por região do espaço de atributos; cada folha deposita a saída da
árvore num lugar próprio; uma transição de agregação (\code{Vote} na floresta,
\code{Decide\_k} no \emph{boosting}) consome todas as saídas e escreve a classe
em \code{Prediction}.

\begin{figure}[ht]
\centering
\begin{tikzpicture}[
  place/.style={circle,draw,minimum size=8mm,inner sep=0pt,font=\scriptsize},
  trans/.style={rectangle,draw,fill=black!6,minimum width=13mm,minimum height=6mm,font=\scriptsize},
  leaf/.style={rectangle,draw,fill=accent!10,minimum width=9mm,minimum height=5mm,font=\tiny},
  >=Stealth]
\node[place] (in) {Input};
\node[trans,right=8mm of in] (dist) {Distribute};
\foreach \i/\y in {0/1.6,1/0,2/-1.6}{
  \node[place] (tin\i) at (4.6,\y) {TIn\textsubscript{\i}};
  \node[leaf] (l\i a) at (6.4,\y+0.38) {folha};
  \node[leaf] (l\i b) at (6.4,\y-0.38) {folha};
  \node[place] (tp\i) at (8.2,\y) {TProb\textsubscript{\i}};
  \draw[->] (dist) -- (tin\i);
  \draw[->] (tin\i) -- (l\i a); \draw[->] (tin\i) -- (l\i b);
  \draw[->] (l\i a) -- (tp\i); \draw[->] (l\i b) -- (tp\i);
  \begin{scope}[on background layer]
    \node[draw=soft,dashed,rounded corners,fit=(l\i a)(l\i b),inner sep=3pt] (box\i) {};
  \end{scope}
}
\node[font=\tiny,soft,above=0pt of box0] {página Tree\_0};
\node[trans,right=8mm of tp1] (vote) {Vote};
\node[place,right=8mm of vote] (pred) {\tiny Prediction};
\draw[->] (in) -- (dist);
\foreach \i in {0,1,2}{\draw[->] (tp\i) -- (vote);}
\draw[->] (vote) -- (pred);
\end{tikzpicture}
\caption{Uma floresta de três árvores como rede (simplificada: cada árvore
real tem uma folha por região, com uma guarda sobre os atributos). As árvores
não compartilham nenhum lugar entre \code{Distribute} e \code{Vote}.}
\label{fig:rede-rf}
\end{figure}

Para uma amostra fixa, cada árvore tem exatamente uma folha habilitada, e as
árvores são independentes entre si. Mesmo assim o grafo de ocorrência registra
\emph{todas as ordens} em que elas podem disparar. A Tabela~\ref{tab:medido}
mostra o que foi medido.

\begin{table}[ht]
\centering\small
\begin{tabular}{lrrrrrrrr}
\toprule
Floresta aleatória, $n$ árvores & 1 & 2 & 3 & 4 & 5 & 6 & 7 & 8 \\
\midrule
Marcações alcançáveis & 4 & 6 & 10 & 18 & 34 & 66 & 130 & 258 \\
\bottomrule
\end{tabular}

\vspace{8pt}
\begin{tabular}{lrrrrr}
\toprule
\emph{Boosting}, 3 classes, $M$ estágios & 1 & 2 & 3 & 4 & 5 \\
\midrule
Marcações alcançáveis & 218 & 731 & 1\,730 & 3\,377 & 5\,834 \\
\bottomrule
\end{tabular}
\caption{Tamanho do grafo de ocorrência, medido. A floresta tem exatamente
$2^n+2$ marcações; o \emph{boosting} binário (um canal) é linear em $M$; o
multiclasse é o produto dos tamanhos dos $K$ canais, que andam em paralelo.}
\label{tab:medido}
\end{table}

A fórmula da floresta é exata e conta a história inteira: depois de
\code{Distribute}, cada marcação corresponde a um \emph{subconjunto} das
árvores que já votaram, e há $2^n$ subconjuntos; somam-se a marcação inicial
e a final. A Figura~\ref{fig:cubo} mostra o caso $n=3$: um cubo.

\begin{figure}[ht]
\centering
\begin{tikzpicture}[
  st/.style={circle,draw,minimum size=7mm,inner sep=0pt,font=\scriptsize},
  red/.style={st,draw=warm,very thick,fill=warm!10},
  >=Stealth, x=1.15cm, y=1.0cm]
\node[st] (m0) at (0,0) {$M_0$};
\node[red] (e) at (1.6,0) {$\emptyset$};
\node[st] (a) at (3.2,1.6) {1};
\node[red] (b) at (3.2,0) {0};
\node[st] (c) at (3.2,-1.6) {2};
\node[st] (ab) at (5.0,1.6) {0,1};
\node[st] (ac) at (5.0,0) {1,2};
\node[red] (bc) at (5.0,-1.6) {0,2};
\node[red] (abc) at (6.8,0) {0,1,2};
\node[red] (f) at (8.4,0) {$M_f$};
\draw[->] (m0) -- node[above,font=\tiny]{Distribute} (e);
\draw[->,soft] (e) -- (a); \draw[->,warm,very thick] (e) -- (b); \draw[->,soft] (e) -- (c);
\draw[->,soft] (a) -- (ab); \draw[->,soft] (a) -- (ac);
\draw[->,soft] (b) -- (ab); \draw[->,warm,very thick] (b) -- (bc);
\draw[->,soft] (c) -- (ac); \draw[->,soft] (c) -- (bc);
\draw[->,soft] (ab) -- (abc); \draw[->,soft] (ac) -- (abc); \draw[->,warm,very thick] (bc) -- (abc);
\draw[->,warm,very thick] (abc) -- node[above,font=\tiny]{Vote} (f);
\end{tikzpicture}
\caption{O grafo de ocorrência de uma floresta de três árvores. Cada nó do
meio é o conjunto de árvores que já votaram. Em destaque, o único caminho que
o grafo reduzido explora: $n+3$ nós em vez de $2^n+2$.}
\label{fig:cubo}
\end{figure}

\begin{nota}
\textbf{A conclusão da medição.} A explosão nessas redes é
\emph{exclusivamente} entrelaçamento (\emph{interleaving}) de componentes
independentes. Nenhuma árvore tem mais de um caminho interno para uma amostra
fixa; o que cresce é o número de ordens em que $n$ coisas independentes podem
acontecer. Uma floresta de 100 árvores tem $2^{100}+2 \approx 1{,}3\times
10^{30}$ marcações --- e todas as ordens levam ao mesmo voto.
\end{nota}

% ============================================================================
\chapter{As técnicas avaliadas}
\label{cap:avaliacao}
% ============================================================================

O capítulo~8 de Jensen \& Kristensen apresenta três métodos de redução
desenvolvidos no contexto de CPN: \emph{sweep-line}, simetria e equivalência.
Ele próprio observa que nenhum método funciona para todo tipo de sistema, que
cada um explora uma característica específica do modelo, e que o método
escolhido precisa (i) explorar uma característica presente no sistema e (ii)
preservar as propriedades que se quer verificar. Os três foram avaliados com
esses dois critérios contra as redes medidas acima.

\section{Sweep-line}

\textbf{Ideia.} Se existe uma \emph{medida de progresso} que nunca diminui ao
longo das ocorrências, o grafo pode ser explorado camada por camada e as
camadas já processadas podem ser apagadas da memória: nenhuma marcação futura
volta a elas. O pico de memória cai para o tamanho de uma camada.

\textbf{Aplica-se?} Sim, e trivialmente: as redes são acíclicas, e o número de
transições já disparadas é uma medida de progresso monotônica.

\textbf{Ajuda?} Pouco. A camada mais larga da floresta é a do meio do cubo,
com $\binom{n}{\lfloor n/2\rfloor}$ marcações --- para $n=8$, o pico medido foi
126 de 258 nós, um ganho de cerca de $2\times$, o mesmo patamar da tabela de
estatísticas do próprio capítulo. Para $n=100$ o pico ainda é da ordem de
$10^{29}$. A técnica reduz memória, não o número de estados visitados.

\textbf{Preserva?} Só propriedades de segurança (``existe marcação alcançável
que satisfaz um predicado''). Perderíamos as propriedades de vivacidade
(\code{AF}), a de recuperabilidade (\code{AG EF}), a de \emph{home marking} (que
precisa dos componentes fortemente conexos) e os contraexemplos, porque a
versão básica apaga justamente os caminhos que eles reproduzem.

\section{Simetria}

\textbf{Ideia.} Se componentes são intercambiáveis (os receptores idênticos do
protocolo do capítulo), marcações que diferem só por uma permutação das
identidades viram um único nó.

\textbf{Aplica-se?} Não. O critério de corretude exige que toda guarda e toda
inscrição de arco seja simétrica sob a permutação. As árvores de uma floresta
têm limiares diferentes; trocar a árvore $i$ pela $j$ não é uma simetria da
rede. O mesmo vale para os canais do \emph{boosting}. Só árvores literalmente
idênticas seriam permutáveis, o que quase não ocorre com \emph{bootstrap}.

\section{Equivalência}

\textbf{Ideia.} Generaliza a simetria: o usuário define diretamente quais
marcações são equivalentes, com a obrigação de que marcações equivalentes
tenham conjuntos equivalentes de elementos habilitados e de sucessores.

\textbf{Aplica-se?} Não foi encontrada relação útil que satisfaça o critério.
Duas marcações com subconjuntos diferentes de árvores já disparadas habilitam
transições diferentes; não são equivalentes. Além disso, a prova de
consistência é manual --- e o \cpn{} já removeu uma vez um modo composicional
justamente por falta de prova de preservação.

\section{O que sobra: redução por ordem parcial}

O mesmo capítulo lembra que há métodos desenvolvidos \emph{fora} do contexto de
CPN e aplicáveis a qualquer sistema de transições. Entre eles está a família
que ataca exatamente o fenômeno medido: a \emph{redução por ordem parcial}
(\emph{partial-order reduction}), que explora a \emph{independência} entre
ocorrências para não enumerar ordens que não importam. Sua forma mais antiga e
mais geral são os \emph{conjuntos teimosos} de Valmari (1988, 1991). A
Tabela~\ref{tab:comparacao} resume a avaliação.

\begin{table}[ht]
\centering\small
\begin{tabular}{>{\raggedright\arraybackslash}p{2.3cm}>{\raggedright\arraybackslash}p{3.0cm}>{\raggedright\arraybackslash}p{3.9cm}>{\raggedright\arraybackslash}p{3.2cm}}
\toprule
Técnica & Explora & Nas redes do \pyra{} & Propriedades \\
\midrule
Sweep-line & progresso monotônico & aplicável; pico $\approx\binom{n}{n/2}$, ganho $\sim2\times$ & só segurança \\
Simetria & componentes permutáveis & \textbf{incorreta}: guardas diferentes & todas as invariantes por simetria \\
Equivalência & relação dada pelo usuário & sem relação válida; prova manual & as que a relação respeita \\
\textbf{Ordem parcial} & \textbf{independência} & \textbf{floresta: $2^n+2 \to n+3$} & dependem das condições (cap.~\ref{cap:preserva}) \\
\bottomrule
\end{tabular}
\caption{Resumo da avaliação.}
\label{tab:comparacao}
\end{table}

% ============================================================================
\chapter{A ideia central: independência}
\label{cap:ideia}
% ============================================================================

\section{O diamante}

Duas ocorrências $a$ e $b$ são \emph{independentes} numa marcação $M$ quando
ambas estão habilitadas, disparar uma não desabilita a outra, e as duas ordens
levam à mesma marcação. No grafo isso aparece como um losango
(Figura~\ref{fig:diamante}).

\begin{figure}[ht]
\centering
\begin{tikzpicture}[st/.style={circle,draw,minimum size=8mm,inner sep=0pt,font=\small},>=Stealth]
\node[st] (m) at (0,0) {$M$};
\node[st] (ma) at (-1.6,-1.3) {$M_a$};
\node[st] (mb) at (1.6,-1.3) {$M_b$};
\node[st] (mab) at (0,-2.6) {$M'$};
\draw[->] (m) -- node[above left,font=\small]{$a$} (ma);
\draw[->] (m) -- node[above right,font=\small]{$b$} (mb);
\draw[->] (ma) -- node[below left,font=\small]{$b$} (mab);
\draw[->] (mb) -- node[below right,font=\small]{$a$} (mab);
\end{tikzpicture}
\caption{$a$ e $b$ independentes: as duas ordens chegam a $M'$. Explorar só
uma delas não perde $M'$; perde apenas $M_b$ (ou $M_a$), uma marcação
intermediária.}
\label{fig:diamante}
\end{figure}

A redução explora só um lado do losango. O que se perde é a marcação
intermediária do outro lado. Toda a teoria que segue responde a uma pergunta:
\emph{quando} é seguro perder marcações intermediárias? A resposta depende do
que se quer verificar. Se a pergunta é ``a rede pode travar?'', marcações
intermediárias nunca importam, porque não são mortas. Se a pergunta é ``a
predição nunca tem dois tokens?'', importam só as que podem violar isso. Se é
``toda execução chega a uma decisão?'', importa a ordem relativa de quem mexe
na predição.

\section{Do par ao conjunto}

Com $n$ árvores independentes, há $n!$ ordens das folhas e $2^n$ marcações
intermediárias. Escolher sempre a mesma árvore primeiro --- a de menor índice
que ainda não votou --- percorre um único caminho de $n$ passos. Mas
``independente'' não pode ser um julgamento caso a caso sobre marcações já
calculadas: a redução tem de decidir \emph{antes} de explorar, olhando só a
marcação atual e a estrutura da rede. É isso que o conjunto teimoso
formaliza: um conjunto de transições, calculado na marcação atual, tal que
disparar primeiro um de seus membros nunca impede nada que outra ordem poderia
fazer.

% ============================================================================
\chapter{Dependência lida dos arcos}
\label{cap:dependencia}
% ============================================================================

O \cpn{} não tenta descobrir a independência executando a rede. Ele a lê da
\emph{estrutura}: de quais lugares cada transição tira tokens, em quais
deposita, e quais apenas consulta. Isso é feito uma vez por rede.

\section{Tirar, dar, testar}

Para cada transição $t$ e cada lugar $p$ ligado a ela:

\begin{itemize}[leftmargin=1.4em]
  \item $t$ \textbf{tira} de $p$ se há um arco de entrada de $p$ para $t$;
  \item $t$ \textbf{dá} a $p$ se há um arco de saída de $t$ para $p$;
  \item $t$ \textbf{testa} $p$ se há um arco de entrada e um de saída
        \emph{com a mesma inscrição}. O token sai e volta igual: o efeito
        líquido é nulo, mas $t$ precisa que ele esteja lá. Nesse caso $p$ não
        conta como ``tirado'' nem ``dado''.
\end{itemize}

Os \emph{lugares de entrada} de $t$ são os que ela tira ou testa; os
\emph{lugares que $t$ muda} são os que ela tira ou dá.

\begin{nota}
\textbf{Por que estrutural e não semântica.} Uma transição que dá a um lugar
pode, para certos valores, não depositar nada (uma inscrição
\code{if \ldots{} then 1`x else empty}). A leitura estrutural a trata como
``pode dar'' sempre. Isso é uma \emph{sobre-aproximação} da dependência: a
redução pode manter mais estados do que precisaria, nunca menos. A correção não
depende de entender o que as inscrições calculam.
\end{nota}

\section{As relações derivadas}

A partir disso, para cada lugar $p$, ficam indexados: quem dá a $p$
(\emph{produtores}), quem tira de $p$, e quem precisa de $p$ (tem $p$ como
lugar de entrada). E para cada transição $t$ fica pronto o seu \emph{conjunto
de conflito}: todas as transições que mudam algum lugar de entrada de $t$, e
todas as que precisam de algum lugar de que $t$ tira. Na floresta, o conjunto
de conflito de uma folha da árvore~$i$ é: as outras folhas da árvore~$i$ (que
tiram de \code{TIn}$_i$ como ela) e \code{Distribute} (que dá a
\code{TIn}$_i$). Nenhuma transição de outra árvore aparece --- e isso é tudo o
que a redução precisa saber para não entrelaçar as árvores.

% ============================================================================
\chapter{O conjunto teimoso de uma marcação}
\label{cap:teimoso}
% ============================================================================

\section{As duas regras de fechamento}

Em cada marcação $M$, o conjunto teimoso $S$ começa com uma ou mais
transições-semente e é fechado por duas regras, aplicadas até nada mudar:

\begin{enumerate}[leftmargin=1.6em]
  \item \textbf{Transição habilitada em $S$.} Entram em $S$ todas as que
        mudam algum lugar de entrada dela, e todas as que precisam de algum
        lugar de que ela tira --- o seu conjunto de conflito.

        \emph{Efeito:} nenhuma transição fora de $S$ consegue desabilitá-la,
        nem ser desabilitada por ela, nem alterar os tokens dos quais suas
        ligações dependem.
  \item \textbf{Transição desabilitada em $S$.} Entram os \emph{habilitadores}
        dela. Se algum lugar de entrada tem menos tokens do que qualquer
        ligação exige, basta escolher \emph{um} desses lugares e incluir seus
        produtores --- o de menos produtores. Se nenhum está curto (a transição
        está desabilitada pelos \emph{valores} dos tokens ou pela guarda),
        entram os produtores de todos os lugares de entrada.

        \emph{Efeito:} nenhuma transição fora de $S$ consegue habilitá-la.
        Tirar tokens nunca habilita nada, por isso só produtores importam.
\end{enumerate}

Os elementos de ligação expandidos em $M$ são então todos os elementos
habilitados das transições de $S$ --- e só eles.

\section{O que as regras garantem}

As regras implementam as duas condições clássicas de Valmari:

\begin{description}[leftmargin=1.4em,font=\normalfont\bfseries]
  \item[D1 (comutação).] Se uma sequência de transições fora de $S$ seguida
        de uma transição $t\in S$ é possível a partir de $M$, então $t$ seguida
        da mesma sequência também é, e chega à mesma marcação. Em outras
        palavras: o que está em $S$ pode sempre ser antecipado.
  \item[D2 (chave).] Existe em $S$ uma transição habilitada que continua
        habilitada qualquer que seja a sequência de transições fora de $S$.
        Pela regra~1, \emph{toda} transição habilitada de $S$ é uma chave.
  \item[D0.] $S$ contém alguma transição habilitada sempre que $M$ não é
        morta. O \cpn{} garante isso escolhendo pelo menos uma semente
        habilitada.
\end{description}

\begin{nota}
\textbf{Uma sutileza de redes coloridas.} Em redes P/T basta proteger a
transição contra ser desabilitada. Em CPN, uma transição habilitada pode ganhar
\emph{novas ligações} quando alguém deposita tokens com outros valores nos seus
lugares de entrada --- e uma ligação nova que só existe depois de uma transição
de fora de $S$ quebraria D1. Por isso a regra~1 inclui em $S$ também quem
\emph{dá} aos lugares de entrada, não apenas quem tira deles. Assim os tokens de
entrada de $t$ ficam intocados por tudo fora de $S$, e as ligações habilitadas
de $t$ são exatamente as mesmas antes e depois.
\end{nota}

\section{Escolhendo entre candidatos}

Há em geral várias sementes possíveis, e cada uma gera um conjunto fechado
diferente. O \cpn{} prefere o candidato que expande o menor número de
elementos de ligação e para assim que encontra um com exatamente um. Como as
duas regras olham uma transição por vez, o fecho de uma união é a união dos
fechos; o fecho de cada transição isolada é calculado uma vez por marcação e
combinado.

\section{Exemplo: a floresta de três árvores, passo a passo}

\begin{enumerate}[leftmargin=1.6em]
  \item Em $M_0$ só \code{Distribute} está habilitada. Ela é expandida (não há
        o que reduzir).
  \item Em $\emptyset$ estão habilitadas a folha certa de cada uma das três
        árvores. Semente: a folha da árvore~0. Regra~1: entram as outras folhas
        da árvore~0 e \code{Distribute}. As outras folhas estão desabilitadas
        (guarda falsa para esta amostra): regra~2, sem lugar curto, entram os
        produtores de \code{TIn}$_0$ --- só \code{Distribute}, que já está.
        \code{Distribute} está desabilitada (\code{Input} vazio): regra~2,
        \code{Input} está curto e não tem produtores. Fim: $S$ tem um único
        elemento habilitado. Só a árvore~0 dispara.
  \item O mesmo acontece em $\{0\}$ com a árvore~1 e em $\{0,1\}$ com a
        árvore~2.
  \item Em $\{0,1,2\}$ só \code{Vote} está habilitada; dispara; a marcação
        final é morta.
\end{enumerate}

O grafo reduzido é o caminho destacado na Figura~\ref{fig:cubo}: $n+3$ nós.
As marcações $\{1\}$, $\{2\}$, $\{0,1\}$\ldots{} nunca são construídas.

% ============================================================================
\chapter{O que se preserva depende da pergunta}
\label{cap:preserva}
% ============================================================================

D0--D2 sozinhas garantem uma coisa: nenhuma marcação morta se perde. Isso
basta para algumas propriedades e não para outras. O \cpn{} constrói até
quatro \emph{tipos} de grafo, cada um com as condições adicionais que a
pergunta exige, e decide cada propriedade no tipo menos explorado que
comprovadamente dá o mesmo veredito do grafo completo.

\section{Grafo \emph{deadlock}: só D0--D2}

\begin{teorema}[Valmari]
Com D0--D2, toda marcação morta alcançável no grafo completo é alcançável no
grafo reduzido, e vice-versa (o reduzido é um subgrafo do completo).
\end{teorema}

\emph{Por quê, informalmente.} Pegue um caminho do grafo completo até uma
marcação morta. Ele precisa conter alguma transição de $S$: senão a chave de
D2 continuaria habilitada até o fim, e a marcação não seria morta. Pela D1, a
primeira transição de $S$ no caminho pode ser antecipada; o reduzido dá esse
passo, e o caminho restante é um passo mais curto. Por indução, o reduzido
chega à mesma marcação morta.

O mesmo argumento mostra mais uma coisa, que o \cpn{} usa:

\begin{teorema}
Com D0--D2, existe uma sequência de ocorrência infinita no grafo completo se e
somente se o grafo reduzido tem um ciclo.
\end{teorema}

\emph{Por quê.} Dada uma sequência infinita, ou ela contém uma transição de
$S$ (que D1 antecipa), ou não contém nenhuma (e então a chave, que nunca é
desabilitada, pode ser disparada antes de qualquer prefixo dela). Em ambos os
casos o reduzido dá um passo e continua havendo uma sequência infinita a
partir dali. Como o grafo é finito, isso fecha um ciclo.

\textbf{Decide:} terminação (\code{AF dead}), ``toda marcação morta tem
predição'' (\code{AG(dead $\to$ pred)}) e, quando a rede termina, o \emph{home
marking}: num grafo finito acíclico os componentes terminais são as marcações
mortas, e existe \emph{home marking} exatamente quando há uma só.

\section{Grafo \emph{safety}: D0--D2 mais um objetivo}

Para invariantes \code{AG $\psi$} com $\psi$ proposicional (``a predição tem no
máximo um token'', ``nunca há duas folhas da mesma árvore habilitadas''), a
pergunta é se alguma marcação que viola $\psi$ é alcançável --- mesmo que não
seja morta. Marcações intermediárias passam a importar, mas só as que podem
violar $\psi$.

A condição extra vem de Kristensen \& Valmari (2000): numa marcação em que
$\psi$ ainda vale, $S$ precisa conter um \emph{up-set} para $\neg\psi$ ---
um conjunto de transições tal que todo caminho que torna $\psi$ falsa dispara
alguma delas.

\begin{teorema}[Kristensen \& Valmari]
Com D0--D2 e o up-set, uma marcação violando $\psi$ é alcançável no grafo
completo se e somente se é alcançável no reduzido. Não é preciso condição de
ciclo.
\end{teorema}

\emph{Por quê.} Pegue o caminho mais curto até uma violação. Ele dispara
alguma transição do up-set, e portanto de $S$; D1 antecipa a primeira; o
caminho restante é estritamente mais curto. O comprimento cai a cada passo ---
não há como ``andar em círculos'' ignorando a violação.

\subsection*{Como o up-set é calculado}

O up-set é calculado na marcação atual, a partir da fórmula e do seu valor
presente. Para cada proposição atômica, a Tabela~\ref{tab:upset} diz quem
consegue mudar seu valor.

\begin{table}[ht]
\centering\small
\begin{tabular}{>{\raggedright\arraybackslash}p{3.1cm}>{\raggedright\arraybackslash}p{4.6cm}>{\raggedright\arraybackslash}p{4.4cm}}
\toprule
Proposição & Hoje verdadeira; para ficar falsa & Hoje falsa; para ficar verdadeira \\
\midrule
algum lugar marcado & tirar de \emph{um} lugar marcado (qualquer um serve) & dar a algum dos lugares \\
todos marcados & tirar de algum dos lugares & dar a \emph{um} lugar vazio \\
$|p|\le k$ para os lugares & dar a algum dos lugares & tirar de \emph{um} lugar acima de $k$ \\
tokens num domínio & dar a algum dos lugares & tirar de \emph{um} lugar com token fora \\
$T$ habilitada & tirar das entradas de \emph{uma} habilitada & habilitadores de cada $t\in T$ \\
$|\mathrm{EN}(T)|\le k$ & habilitadores das desabilitadas, produtores das entradas das habilitadas & tirar das entradas de $T$ \\
\bottomrule
\end{tabular}
\caption{Quem pode virar cada proposição. ``\emph{Um}'' indica que basta
escolher um elemento (o de menor conjunto), porque todos precisam mudar.}
\label{tab:upset}
\end{table}

Os conectivos combinam isso pela lógica de ``quem precisa mudar'': uma
conjunção verdadeira cai se \emph{qualquer} conjunto cair, então os up-sets
dos conjuntos se somam; uma conjunção falsa só fica verdadeira se \emph{todo}
conjunto falso virar, então basta o up-set de um deles, o menor. A disjunção é
o dual, e a negação apenas inverte o valor.

\begin{exemplo}
Propriedade D1 da floresta: \code{AG(EN(Vote) $\to$ todas as árvores
votaram)}. No início, \code{EN(Vote)} é falsa, e a implicação é verdadeira
por causa dela. Para violá-la, \code{Vote} precisa ficar habilitada; o lugar
de entrada curto de menor número de produtores é a saída de uma árvore; o
up-set são as folhas daquela árvore. A busca segue linear.
\end{exemplo}

\section{Grafo \emph{stutter}: visibilidade, singleton e ciclos}

Para \code{AF}, \code{EF}, \code{EU} e combinações como \code{AG EF}, a
\emph{ordem} relativa dos eventos observados importa. A teoria aplicada é a de
Gerth, Kuiper, Peled \& Penczek (1999), que acrescenta a D0--D2:

\begin{description}[leftmargin=1.4em,font=\normalfont\bfseries]
  \item[C2 (visibilidade).] Uma transição é \emph{visível} se muda algum lugar
        que a fórmula observa (os lugares das proposições, ou os lugares de
        entrada das transições cuja habilitação a fórmula consulta). Um
        conjunto reduzido só pode conter transições invisíveis.
  \item[C3 (ciclo).] Se o passo reduzido levaria a uma marcação que está na
        pilha da busca em profundidade (fecharia um ciclo), a marcação é
        expandida por completo. Sem isso, um ciclo poderia adiar para sempre
        uma transição que ninguém escolhe (o ``problema da ignorância'').
  \item[C4 (singleton).] O conjunto reduzido tem \emph{um único} elemento de
        ligação; caso contrário, expande-se tudo.
\end{description}

\begin{teorema}[Gerth, Kuiper, Peled \& Penczek]
Com essas condições o grafo reduzido é bissimilar por gagueira
(\emph{stuttering bisimilar}) ao completo em relação às proposições
observadas. Portanto toda fórmula CTL sem o operador ``próximo'' (CTL-X) tem o
mesmo valor nos dois.
\end{teorema}

\emph{Gagueira} significa: as execuções do reduzido e do completo mostram as
mesmas mudanças das proposições observadas, na mesma ordem, diferindo apenas
em quantas vezes cada valor se repete. Como CTL-X não sabe contar passos,
não distingue as duas. Nenhuma propriedade do catálogo usa ``próximo''.

Por causa de C3, este tipo é explorado em profundidade; os outros dois, em
largura, o que dá contraexemplos mais curtos.

\section{Grafo \emph{full}}

O grafo de ocorrência comum. Usado para as fórmulas que nenhum dos teoremas
cobre --- em particular as que mencionam marcações mortas de um jeito diferente
dos formatos reconhecidos (por exemplo \code{AG EF dead}), e o \emph{home
marking} de uma rede que não termina.

\section{Por que \emph{dead} é tratado à parte}

``A marcação é morta'' não é uma proposição sobre lugares: \emph{qualquer}
transição pode torná-la verdadeira ao ser a última. Tratá-la como observada
tornaria toda transição visível e anularia a redução. Por isso ela só aparece
nos formatos que o grafo \emph{deadlock} decide sem visibilidade:
\code{AF dead}, \code{AG(dead $\to\psi$)}, \code{EF(dead $\wedge\chi$)} e a
negação deste, com $\psi$, $\chi$ sem \emph{dead}.

% ============================================================================
\chapter{Da propriedade ao grafo}
\label{cap:pipeline}
% ============================================================================

\begin{figure}[ht]
\centering
\begin{tikzpicture}[
  box/.style={rectangle,draw,rounded corners=2pt,fill=callout,minimum height=8mm,
              text width=2.7cm,align=center,font=\small},
  kind/.style={rectangle,draw,rounded corners=2pt,minimum height=7mm,
               text width=1.8cm,align=center,font=\scriptsize},
  >=Stealth]
\node[box] (p) at (0,0) {propriedade (fórmula)};
\node[box] (s) at (4,0) {divisão em partes};
\node[box] (c) at (8,0) {classificação de cada parte};
\node[kind,fill=accent!8] (k1) at (1.3,-1.8) {deadlock};
\node[kind,fill=accent!8] (k2) at (3.5,-1.8) {safety};
\node[kind,fill=accent!8] (k3) at (5.7,-1.8) {stutter};
\node[kind,fill=black!6] (k4) at (7.9,-1.8) {full};
\node[box] (g) at (2.4,-3.6) {grafos reduzidos (com cache)};
\node[box] (v) at (6.8,-3.6) {veredito = conjunção das partes};
\draw[->] (p) -- (s); \draw[->] (s) -- (c);
\foreach \k in {k1,k2,k3,k4}{\draw[->] (c.south) -- (\k.north); \draw[->] (\k.south) -- (g.north);}
\draw[->] (g) -- (v);
\end{tikzpicture}
\caption{Cada propriedade é dividida, cada parte é classificada, e cada
classe é decidida no seu grafo.}
\label{fig:pipeline}
\end{figure}

\section{Divisão em partes}

Uma parte que observa menos lugares tem menos transições visíveis e um up-set
menor, e portanto reduz mais. Antes de classificar, toda propriedade é
dividida em partes independentes cuja conjunção é ela mesma:

\begin{itemize}[leftmargin=1.4em]
  \item \textbf{quantificador ``para toda árvore/canal''}: vira uma parte por
        árvore ou canal, cada uma com o nome do membro, que aparece no
        contraexemplo;
  \item \textbf{\code{AG} de conjunção}: \code{AG}$(a\wedge b) =$
        \code{AG}$\,a\,\wedge$ \code{AG}$\,b$;
  \item \textbf{limite sobre muitos lugares}: ``todo lugar tem no máximo um
        token'' (A7) vira uma invariante por lugar. Sem isso, o up-set seria
        ``todos os produtores de todos os lugares'' --- isto é, tudo --- e nada
        seria reduzido.
\end{itemize}

\section{Classificação}

Cada parte é examinada pela forma da fórmula: \code{AF dead} e as formas com
\emph{dead} reconhecidas vão para \emph{deadlock}; \code{AG}, \code{!EF} e
\code{EF} de fórmulas proposicionais sem \emph{dead} vão para \emph{safety}
(com o $\psi$ cuja violação é o objetivo); o resto sem \emph{dead} vai para
\emph{stutter} (com o seu conjunto de transições visíveis); o que sobrar vai
para \emph{full}. A Tabela~\ref{tab:catalogo} mostra o resultado para o
catálogo das redes de árvores.

\begin{table}[ht]
\centering\small
\begin{tabular}{llll}
\toprule
Id & Propriedade & Grafo & Partes (RF, 100 árvores) \\
\midrule
A1 & terminação & deadlock & 1 \\
A2 & sem travamento espúrio & deadlock & 1 \\
A3 & decisão inevitável (\code{AF pred}) & stutter & 1 \\
A4 & recuperabilidade (\code{AG EF pred}) & stutter & 1 \\
A5 & saída única & safety & 1 \\
A6 & rótulo válido & safety & 1 \\
A7 & \emph{safeness} & safety & 202 (uma por lugar) \\
A8 & \emph{home marking} & deadlock\textsuperscript{*} & --- \\
B1 & determinismo das folhas & safety & 100 \\
B2 & seleção inevitável de folha & stutter & 100 \\
B3 & saída única por árvore & safety & 100 \\
C1, C2 & precedência dos estágios (GBDT) & stutter & por canal \\
C3a & decisão inevitável (GBDT) & stutter & 1 \\
C3b & decisão determinística (GBDT) & safety & 1 \\
D1 & votação completa & safety & 1 \\
D2a & agregação inevitável (\code{AF EN(Vote)}) & stutter & 1 (\textbf{não reduz}) \\
D2b & agregação única & safety & 1 \\
\bottomrule
\end{tabular}
\caption{Onde cada propriedade do catálogo é decidida. \textsuperscript{*}Se a
rede termina; senão, o grafo completo. B4 (estrutural) e PC
(teste de conformidade) não usam grafo e não mudam.}
\label{tab:catalogo}
\end{table}

\section{Caches}

Grafos são construídos sob demanda e guardados por tipo e parâmetro: um único
grafo \emph{deadlock} serve A1, A2 e A8; grafos \emph{stutter} com o mesmo
conjunto visível são compartilhados. Além disso, as centenas de grafos
\emph{safety} de uma mesma entrada percorrem quase sempre as mesmas primeiras
marcações em ordens diferentes, então os elementos habilitados de cada
marcação são lembrados entre grafos. E dentro de um grafo, ao passar de uma
marcação para a seguinte, só são reavaliadas as transições que precisam de
algum lugar que a ocorrência mudou --- as outras mantêm exatamente as mesmas
ligações.

% ============================================================================
\chapter{Detalhes que importam para a correção}
\label{cap:detalhes}
% ============================================================================

\section{Habilitação num grafo reduzido}

Propriedades como ``nunca há duas folhas habilitadas'' ou ``\code{Vote} acaba
habilitada'' consultam a \emph{habilitação} numa marcação. No grafo completo,
habilitação e arcos de saída coincidem. No reduzido, não: os arcos são só os
elementos expandidos. Por isso cada nó de um grafo reduzido guarda a lista
completa de elementos habilitados encontrada durante a busca, e as
proposições de habilitação são avaliadas sobre ela, nunca sobre os arcos. Sem
isso, uma redução esconderia exatamente o não-determinismo que B1 procura.

\section{O não-determinismo continua visível}

Se uma rede mutante tem duas folhas da mesma árvore habilitadas ao mesmo tempo
(uma sobreposição de guardas alcançada pela amostra), as duas tiram do mesmo
lugar; pela regra~1 entram juntas em qualquer conjunto teimoso. A redução
nunca escolhe uma e esconde a outra. Isso aparece nos testes: os mutantes de
sobreposição continuam reprovados com a redução ativa.

\section{Contraexemplos}

Todo grafo reduzido é um subgrafo do completo, então todo caminho nele é uma
sequência de ocorrência real da rede. Os contraexemplos encontrados num grafo
reduzido são construídos e certificados exatamente como antes: o
\emph{replay}, que redispara a sequência com a regra de disparo da rede e
confere a marcação final, continua sendo a prova de que não foram inventados.
Eles podem não ser os mais curtos do grafo completo, apenas válidos.

\section{Orçamento}

Cada grafo tem o mesmo limite de marcações que o modo padrão. Uma parte cujo
grafo estoura o limite fica \emph{sem veredito}; a propriedade inteira fica
sem veredito se nenhuma outra parte falhou, e o relatório diz qual grafo
estourou e por quê. Um veredito nunca é chutado.

\section{O relatório}

No modo reduzido, o relatório mostra quantos grafos foram construídos, de
quais tipos, quantos nós no total e o maior; o tamanho do grafo
\emph{deadlock}; e, abaixo de cada propriedade, em qual tipo de grafo ela foi
decidida (por exemplo ``safety x202''). Como o \cpntools{} não tem essa redução,
os números esperados do oráculo (nós, arcos, componentes) não são produzidos
nesse modo: a comparação com o \cpntools{} continua sendo feita no modo
padrão.

% ============================================================================
\chapter{O limite: visibilidade}
\label{cap:limite}
% ============================================================================

Uma propriedade do catálogo não reduz: D2a, \code{AF EN(Vote)}, ``a agregação
acaba habilitada''. A habilitação de \code{Vote} depende dos lugares de saída
de \emph{todas} as árvores, e toda folha dá a um deles. Logo toda folha é
visível, nenhum conjunto reduzido satisfaz C2, e o grafo \emph{stutter} desta
propriedade é o grafo completo. Até 17 árvores ($2^{17}+2 = 131\,074$ marcações) isso ainda cabe no orçamento
padrão; a partir daí D2a fica sem veredito, com a explicação ``$k$ de $m$
transições mudam o que a propriedade observa e não podem ser reordenadas''.

Isso não é um defeito da implementação: é o limite conhecido da redução por
ordem parcial para propriedades de vivacidade. A condição estática de
visibilidade não sabe que só a \emph{última} folha muda o valor da
proposição --- isso depende de quais árvores já votaram, informação que só a
exploração tem. Duas observações tornam o limite pouco grave na prática:

\begin{itemize}[leftmargin=1.4em]
  \item A3 (\code{AF pred}, ``toda execução chega a uma decisão'') observa só
        \code{Prediction}, é decidida em grafo de $n+3$ nós e já garante o que
        D2a garante para uma rede que produz predição apenas via \code{Vote}.
  \item D1 (``\code{Vote} só habilita quando todas votaram''), A1 e A2 cobrem o
        lado de segurança da agregação, todos reduzidos.
\end{itemize}

\begin{nota}
\textbf{Possível continuação.} O capítulo~8 observa que métodos de redução podem
ser combinados. Para D2a, a combinação natural é com o \emph{sweep-line}: o
grafo completo dessa propriedade é acíclico e tem medida de progresso trivial,
então poderia ser percorrido camada a camada com memória proporcional à camada
mais larga. Isso troca o limite de memória por tempo, sem mudar o número de
estados; não foi implementado.
\end{nota}

% ============================================================================
\chapter{Validação}
\label{cap:validacao}
% ============================================================================

Uma redução só vale se for invisível nas respostas. O critério adotado foi o
mais forte disponível sem um assistente de provas: comparar veredito por
veredito com o grafo completo, em redes construídas para falhar cada
propriedade.

\section{Concordância de vereditos}

\begin{itemize}[leftmargin=1.4em]
  \item \textbf{No \cpn{}}: as duas redes de referência (\emph{boosting}
        binário e floresta multiclasse) e \emph{todos} os seus mutantes --- um
        por propriedade, incluindo ciclo, travamento, saída dupla, rótulo
        inválido, sobreposição de guardas, votação parcial, estágio pulado ---
        somando 28 redes e \textbf{435 vereditos}: \textbf{nenhuma
        divergência}.
  \item \textbf{No \pyra{}}: árvore de decisão, floresta, \emph{boosting}
        binário e de três classes, nas versões inicial e refinada e em todos
        os mutantes de cada uma: mesma concordância total, e o caminho completo
        (\emph{model\_check} com redução) aprova as redes corretas.
  \item \textbf{Escala}: florestas de até 16 árvores e \emph{boosting} de três
        classes com até 5 estágios, comparados com o grafo completo: nenhuma
        divergência.
\end{itemize}

\section{Propriedades dos próprios grafos}

\begin{itemize}[leftmargin=1.4em]
  \item todo arco de um grafo reduzido é um arco do grafo completo;
  \item as marcações mortas do grafo \emph{deadlock} são exatamente as do
        completo, em todas as redes e mutantes;
  \item a floresta de $n$ árvores tem $2^n+2$ marcações no completo e $n+3$ no
        reduzido;
  \item num grafo \emph{stutter}, toda marcação expandida parcialmente
        expandiu exatamente um elemento, invisível;
  \item os contraexemplos encontrados com a redução reexecutam na rede.
\end{itemize}

\section{O modo padrão não mudou}

Os testes de concordância com o \cpntools{} 4.0.1 (executado sem interface,
via Access/CPN) continuam passando no modo padrão: a redução é um caminho
paralelo, ativado só quando pedido.

% ============================================================================
\chapter{Resultados}
\label{cap:resultados}
% ============================================================================

\begin{table}[ht]
\centering\small
\setlength{\tabcolsep}{4.5pt}
\begin{tabular}{lrrrrrr}
\toprule
& \multicolumn{2}{c}{Grafo completo} & \multicolumn{4}{c}{Com redução} \\
\cmidrule(lr){2-3}\cmidrule(lr){4-7}
Modelo & nós & tempo & \emph{deadlock} & maior grafo & tempo & decididas \\
\midrule
RF, 2 árvores & 6 & 0{,}00\,s & 5 & 6 & 0{,}00\,s & todas \\
RF, 4 & 18 & 0{,}00\,s & 7 & 18 & 0{,}01\,s & todas \\
RF, 8 & 258 & 0{,}06\,s & 11 & 258 & 0{,}03\,s & todas \\
RF, 12 & 4\,098 & 1{,}39\,s & 15 & 4\,098 & 0{,}28\,s & todas \\
RF, 16 & 65\,538 & 36{,}2\,s & 19 & 65\,538 & 5{,}4\,s & todas \\
RF, 20 & $\approx10^{6}$ & --- & 23 & (orçamento) & 19\,s & menos D2a \\
RF, 50 & $\approx10^{15}$ & --- & 53 & (orçamento) & 37\,s & menos D2a \\
RF, 100 & $\approx10^{30}$ & --- & 103 & (orçamento) & 70\,s & menos D2a \\
\midrule
GBDT 3 cl., 1 estágio & 218 & 0{,}01\,s & 18 & 22 & 0{,}03\,s & todas \\
GBDT 3 cl., 3 & 1\,730 & 0{,}17\,s & 36 & 40 & 0{,}15\,s & todas \\
GBDT 3 cl., 5 & 5\,834 & 0{,}78\,s & 54 & 58 & 0{,}48\,s & todas \\
\bottomrule
\end{tabular}
\caption{Uma amostra por rede, orçamento de 200\,000 marcações por grafo. Nas
florestas o ``maior grafo'' é sempre o de D2a, que não reduz; ele também
domina o tempo. A6 não entra na contagem das florestas porque o domínio de
rótulos não foi declarado nesse experimento.}
\label{tab:resultados}
\end{table}

Na floresta de 100 árvores, tirando D2a, \textbf{todas as outras 12
propriedades de grafo} (B4 é estrutural e não usa grafo) são decididas em
cerca de 6\,s, construindo 507 grafos
reduzidos de no máximo 103 nós cada; A7 sozinha usa 202 grafos (um por
lugar) e leva cerca de 2\,s. Os 64\,s restantes são D2a esgotando o orçamento.
No \emph{boosting} multiclasse, onde nenhuma propriedade observa todas as
árvores, tudo reduz: o maior grafo de qualquer propriedade tem 58 nós contra
5\,834 do completo.

\begin{figure}[H]
\centering
\begin{tikzpicture}
\begin{semilogyaxis}[
  width=0.82\linewidth, height=6.2cm,
  xlabel={árvores na floresta ($n$)}, ylabel={marcações},
  xmin=0, xmax=101, ymin=1, ymax=1e31,
  legend pos=north west, legend style={font=\small,draw=none,fill=none},
  grid=major, grid style={black!8}, tick label style={font=\small}]
\addplot[warm,thick,domain=1:100,samples=100] {2^x+2};
\addlegendentry{grafo completo, $2^n+2$}
\addplot[soft,thick,dashed,domain=1:100,samples=100] {2^x/sqrt(1.5708*x)};
\addlegendentry{pico do sweep-line, $\approx\binom{n}{n/2}$}
\addplot[accent,very thick,domain=1:100,samples=100] {x+3};
\addlegendentry{grafo \emph{deadlock} reduzido, $n+3$}
\addplot[only marks,mark=*,mark size=1.6pt,warm] coordinates {(2,6)(4,18)(8,258)(12,4098)(16,65538)};
\addplot[only marks,mark=*,mark size=1.6pt,accent] coordinates {(2,5)(4,7)(8,11)(12,15)(16,19)(20,23)(50,53)(100,103)};
\end{semilogyaxis}
\end{tikzpicture}
\caption{Crescimento em escala logarítmica. Pontos: medidos. O sweep-line
acompanha o grafo completo a uma distância constante na escala log; a redução
por ordem parcial muda a classe de crescimento.}
\label{fig:crescimento}
\end{figure}

% ============================================================================
\chapter{Uso e limitações}
\label{cap:uso}
% ============================================================================

\section{Como ativar}

A redução é opcional e o modo padrão não mudou. Ela é ativada por um
parâmetro, \code{reduction}, com valor \code{"stubborn"}, aceito pela
verificação do \cpn{} (na API e na linha de comando, como opção
\code{-{}-reduction stubborn}) e repassado pelo \pyra{} na verificação do modelo
e no \emph{pipeline} verificado. Nenhuma propriedade, perfil ou rede precisa
ser alterada: a divisão e a classificação são automáticas e valem para
qualquer perfil, não só o das redes de árvores.

\section{Quando usar}

\begin{itemize}[leftmargin=1.4em]
  \item \textbf{Florestas com mais de uns 12 árvores e \emph{boosting}
        multiclasse}: é onde a explosão está, e onde a redução muda a classe
        de crescimento.
  \item \textbf{Comparação com o \cpntools{}}: use o modo padrão; o oráculo
        compara também o tamanho do grafo, que a redução muda de propósito.
  \item \textbf{Redes pequenas}: tanto faz; a redução constrói vários grafos
        pequenos em vez de um, e pode não ser mais rápida.
\end{itemize}

\section{Limitações}

\begin{itemize}[leftmargin=1.4em]
  \item \textbf{Visibilidade} (capítulo~\ref{cap:limite}): propriedades de
        vivacidade que observam lugares mexidos por todos os componentes não
        reduzem.
  \item \textbf{Dependência estrutural}: a leitura dos arcos sobre-aproxima a
        dependência real; uma rede cujos componentes compartilham lugares
        ``por precaução'' reduz menos. Nunca reduz errado por isso.
  \item \textbf{Fórmulas com \emph{dead}} fora dos formatos reconhecidos vão
        para o grafo completo.
  \item \textbf{Por entrada}: como no modo padrão, cada amostra na marcação
        inicial tem seus próprios grafos; a redução não generaliza vereditos
        de uma amostra para outra. Para ``toda entrada'', a análise estrutural
        das guardas (B4) continua sendo a ferramenta.
\end{itemize}

% ============================================================================
\chapter*{Referências}
\addcontentsline{toc}{chapter}{Referências}
% ============================================================================

\begin{itemize}[leftmargin=1.4em,itemsep=3pt]
\small
\item C.~Baier, J.-P.~Katoen. \emph{Principles of Model Checking}. MIT Press,
      2008. Capítulo~8 (redução por ordem parcial).
\item E.~M.~Clarke, O.~Grumberg, D.~Kroening, D.~Peled, H.~Veith.
      \emph{Model Checking}, 2.ª ed. MIT Press, 2018.
\item R.~Gerth, R.~Kuiper, D.~Peled, W.~Penczek. A partial order approach to
      branching time logic model checking. \emph{Information and Computation}
      150(2), 1999.
\item P.~Godefroid. \emph{Partial-Order Methods for the Verification of
      Concurrent Systems}. LNCS 1032, Springer, 1996.
\item K.~Jensen, L.~M.~Kristensen. \emph{Coloured Petri Nets: Modelling and
      Validation of Concurrent Systems}. Springer, 2009. Capítulo~8
      (\emph{Advanced State Space Methods}).
\item L.~M.~Kristensen, A.~Valmari. Improved question-guided stubborn set
      methods for state properties. \emph{ICATPN 2000}, LNCS 1825, 2000.
\item D.~Peled. All from one, one for all: on model checking using
      representatives. \emph{CAV 1993}, LNCS 697, 1993.
\item A.~Valmari. Stubborn sets for reduced state space generation.
      \emph{Advances in Petri Nets 1990}, LNCS 483, 1991.
\item A.~Valmari, H.~Hansen. Stubborn set intuition explained.
      \emph{Transactions on Petri Nets and Other Models of Concurrency XII},
      LNCS 10470, 2017.
\end{itemize}

\end{document}
